\SourcePage{GELFAND-P000468}{unlabelled}
\underline{13/\mbox{$\overline{\underline{\mathrm{XI}}}$}-72.}
\quad\underline{Gelfand sem..}

\underline{L. Vaserstein.}

\uuline{\mbox{$A$}} — assoc. ring (not necessarily with $1$).

\uuline{\mbox{$iA$}} — ring with a formally adjoined unit.

$iA/A=\mathbb Z$,\quad $A$ — two-sided ideal in $iA$.

$GA$ — mult. group of the ring.

$GA\subset A$

$(a,b)\longrightarrow a+b+ab$\qquad This is a binary assoc. oper.
\[
\Vert\qquad (e+a)(e+b)-e=a+b+ab
\]
\sout{\mbox{$iA$}}\quad \uuline{\mbox{$GA$}} — set of invertible elements with respect to this operation.
\Unclear{GELFAND-P000468}{initial-cancel}{The short crossed-out entry before $GA$ may read $iA$.}
\end{GelfandAdjacentEvidence}
\begin{GelfandAdjacentEvidence}{GELFAND-P000468}{2}


\[
a\longrightarrow1+a,\qquad A\longrightarrow1+A
\]
\uwave{homomorphism}.\quad
\begin{tabular}[t]{@{}l@{}}
\underline{into the ordinary}\\
\underline{invert. elems}.
\end{tabular}

\[
\begin{gathered}
f:A\longrightarrow B\\
f(GA)\subset GB
\end{gathered}
\]
\Unclear{GELFAND-P000468}{image-containment}{The sign before $GB$ is overwritten. It is provisionally set as $\subset$.}
\end{GelfandAdjacentEvidence}
\begin{GelfandAdjacentEvidence}{GELFAND-P000468}{3}


$i,G$ — functors to the cat. of rings with $1$ and the cat. of groups.

\uuline{\mbox{$M_kA$}} — matrices $k\times k$ with ordinary multiplication.

\uuline{\mbox{$\widetilde\Omega A$}}$=\{a=(a_{ij})\}$,
in $\forall$ row and column a finite no. of elems. $\ne0$,
(this is a ring)

\uuline{\mbox{$MA$}} — subring of finite matrices (a finite number of elems. $\ne0$).
(two-sided ideal in $\widetilde\Omega A$).

\uuline{\mbox{$\Omega A$}}$=\widetilde\Omega A/MA$
\Unclear{GELFAND-P000468}{omega-letter}{The letter before $A$ on the left is overwritten. It reads $\Omega$.}
\end{GelfandAdjacentEvidence}
\begin{GelfandAdjacentEvidence}{GELFAND-P000468}{4}


$M,\widetilde\Omega,\Omega$ — functors.

\[
\underset{\vert}{\text{\uuline{\mbox{$\Theta A$}}}}\subset MA
\]
upper triangular $a=(a_{ij})$:\quad $a_{ij}=0$ when $i\le j$ \sout{\ldots}
\SourceFigure{GELFAND-P000468}{theta-link}{A short vertical stroke is drawn below the double underline under $\Theta A$.}
\end{GelfandAdjacentEvidence}
\begin{GelfandAdjacentEvidence}{GELFAND-P000468}{5}

\Unclear{GELFAND-P000468}{triangular-convention}{The manuscript has “upper triangular” and $a_{ij}=0$ when $i\le j$.}
\end{GelfandAdjacentEvidence}
\begin{GelfandAdjacentEvidence}{GELFAND-P000468}{6}

\Incomplete{GELFAND-P000468}{cancel-inequality}{Several signs one above another are crossed out after the inequality.}
\end{GelfandAdjacentEvidence}
\begin{GelfandAdjacentEvidence}{GELFAND-P000468}{7}


\underline{Remark}. If $N\xrightarrow{\sim}N$ is a bijection, there is a map of the matrix ring
(\sout{th.} autom.) — simultaneous perm. of rows and columns.
\Unclear{GELFAND-P000468}{cancel-auto}{The crossed-out abbreviation before “autom.” may read «эт».}
\end{GelfandAdjacentEvidence}
\begin{GelfandAdjacentEvidence}{GELFAND-P000468}{8}

(inner autom. perm. of rows and columns)
\Unclear{GELFAND-P000468}{inner-ambient}{The parenthesis reads «внутр. автом.» (“inner autom.”).}
\end{GelfandAdjacentEvidence}
\begin{GelfandAdjacentEvidence}{GELFAND-P000468}{9}


Even if $N\longrightarrow N$ is any map, it is assigned an \underline{endomorphism}.
\Unclear{GELFAND-P000468}{next-word}{The first word reads «Даже» (“even”).}
\end{GelfandAdjacentEvidence}
\begin{GelfandAdjacentEvidence}{GELFAND-P000468}{10}

\Unclear{GELFAND-P000468}{arbitrary-reindexing}{The last word is underlined and reads “endomorphism.”}
\end{GelfandAdjacentEvidence}
\begin{GelfandAdjacentEvidence}{GELFAND-P000468}{11}


$GMA$ — mult. group, $MA$,\quad $GMA\subset MA$.

\uwave{Simplicial structure in \mbox{$MA$}.}

$X\subset GMA$ is called a simplex, if by a left shift and conjugation by a permutation matrix it can be taken into $\Theta A$
\Unclear{GELFAND-P000468}{left-shift-convention}{The manuscript reads «левым сдвигом» (“by a left shift”).}
\end{GelfandAdjacentEvidence}
\begin{GelfandAdjacentEvidence}{GELFAND-P000468}{12}


Note that $G\Theta A=\Theta A$ as a set (they are nilpotent)
\Unclear{GELFAND-P000468}{final-word}{The last letters are overwritten. The word reads “nilpotent.”}
\end{GelfandAdjacentEvidence}



\Needspace{7\baselineskip}
\begin{GelfandAdjacentEvidence}{GELFAND-P000469}{1}
